% ============================================================
%  PRESENTACIÓN BEAMER – HPC: Potencial de Morse
%  Camilo Huertas · Sebastián Rodriguez
%  Universidad Distrital Francisco José de Caldas
% ============================================================
\documentclass[aspectratio=169,10pt]{beamer}

% ─── TEMA BASE ───────────────────────────────────────────────
\usetheme{Madrid}
\usecolortheme{default}

% ─── PALETA DE COLORES PERSONALIZADA ────────────────────────
\definecolor{UDprimary}{RGB}{0,60,120}
\definecolor{UDsecond}{RGB}{220,50,50}
\definecolor{UDlight}{RGB}{230,240,255}
\definecolor{UDdark}{RGB}{20,30,60}
\definecolor{codegreen}{rgb}{0,0.55,0.2}
\definecolor{codegray}{rgb}{0.5,0.5,0.5}
\definecolor{codepurple}{rgb}{0.50,0,0.70}
\definecolor{backcolour}{RGB}{245,245,240}
\definecolor{UDaccent}{RGB}{255,180,0}
\definecolor{UDsuccess}{RGB}{30,140,60}

% ─── APLICAR COLORES AL TEMA ─────────────────────────────────
\setbeamercolor{palette primary}{bg=UDprimary,fg=white}
\setbeamercolor{palette secondary}{bg=UDdark,fg=white}
\setbeamercolor{palette tertiary}{bg=UDsecond,fg=white}
\setbeamercolor{palette quaternary}{bg=UDprimary!80,fg=white}
\setbeamercolor{structure}{fg=UDprimary}
\setbeamercolor{section in toc}{fg=UDprimary}
\setbeamercolor{block title}{bg=UDprimary,fg=white}
\setbeamercolor{block body}{bg=UDlight}
\setbeamercolor{block title alerted}{bg=UDsecond,fg=white}
\setbeamercolor{block body alerted}{bg=UDsecond!10}
\setbeamercolor{block title example}{bg=UDprimary!70,fg=white}
\setbeamercolor{block body example}{bg=UDprimary!8}
\setbeamercolor{frametitle}{bg=UDprimary,fg=white}
\setbeamercolor{title}{fg=white}
\setbeamercolor{author}{fg=UDlight}
\setbeamercolor{date}{fg=UDlight}
\setbeamercolor{itemize item}{fg=UDsecond}
\setbeamercolor{itemize subitem}{fg=UDprimary}
\setbeamercolor{enumerate item}{fg=UDsecond}

% ─── FUENTES ─────────────────────────────────────────────────
\usepackage[T1]{fontenc}
\usepackage[utf8]{inputenc}
\usepackage[spanish]{babel}
\usepackage{lmodern}
\usepackage{microtype}

% ─── PAQUETES ESENCIALES ─────────────────────────────────────
\usepackage{amsmath,amssymb}
\usepackage{braket}
\usepackage{graphicx}
\usepackage{booktabs}
\usepackage{tabularx}
\usepackage{multirow}
\usepackage{listings}
\usepackage{xcolor}
\usepackage{tikz}
\usetikzlibrary{shapes.geometric,arrows.meta,positioning,shadows,fit,decorations.pathreplacing}
\usepackage{calc}
\usepackage{pgfplots}
\pgfplotsset{compat=1.18}
\usepackage{tcolorbox}
\tcbuselibrary{skins,breakable}
\usepackage{fontawesome5}

% ─── ESTILO DE CÓDIGO ────────────────────────────────────────
\lstdefinestyle{cppBeamer}{
  backgroundcolor=\color{backcolour},
  commentstyle=\color{codegreen}\itshape,
  keywordstyle=\color{UDprimary}\bfseries,
  keywordstyle=[2]\color{UDsecond}\bfseries,
  numberstyle=\tiny\color{codegray},
  stringstyle=\color{codepurple},
  basicstyle=\ttfamily\tiny,
  breakatwhitespace=false,
  breaklines=true,
  captionpos=b,
  keepspaces=true,
  numbers=left,
  numbersep=4pt,
  showspaces=false,
  showstringspaces=false,
  tabsize=2,
  frame=single,
  rulecolor=\color{UDprimary!40},
  xleftmargin=8pt,
  xrightmargin=4pt
}
\lstdefinestyle{fortranBeamer}{
  backgroundcolor=\color{backcolour},
  commentstyle=\color{codegreen}\itshape,
  keywordstyle=\color{UDsecond}\bfseries,
  numberstyle=\tiny\color{codegray},
  stringstyle=\color{codepurple},
  basicstyle=\ttfamily\tiny,
  breaklines=true,
  numbers=left,
  numbersep=4pt,
  frame=single,
  rulecolor=\color{UDsecond!40},
  xleftmargin=8pt,
  xrightmargin=4pt,
  language=Fortran
}
\lstset{
  style=cppBeamer,
  inputencoding=utf8,
  extendedchars=true,
  literate=
    {á}{{\'a}}1 {é}{{\'e}}1 {í}{{\'i}}1 {ó}{{\'o}}1 {ú}{{\'u}}1
    {Á}{{\'A}}1 {É}{{\'E}}1 {Í}{{\'I}}1 {Ó}{{\'O}}1 {Ú}{{\'U}}1
    {ñ}{{\~n}}1 {Ñ}{{\~N}}1 {ü}{{\"u}}1 {Ü}{{\"U}}1
    {—}{{--}}1 {–}{{--}}1
}

% ─── TCOLORBOX SHORTCUTS ─────────────────────────────────────
\newtcolorbox{conclusionbox}[1]{
  colback=UDlight,colframe=UDprimary,
  fonttitle=\bfseries\small,title=#1,
  left=4pt,right=4pt,top=3pt,bottom=3pt
}
\newtcolorbox{alertbox}[1]{
  colback=UDsecond!10,colframe=UDsecond,
  fonttitle=\bfseries\small,title=#1,
  left=4pt,right=4pt,top=3pt,bottom=3pt
}
\newtcolorbox{physbox}[1]{
  colback=UDprimary!8,colframe=UDprimary!60,
  fonttitle=\bfseries\small,title=#1,
  left=4pt,right=4pt,top=3pt,bottom=3pt
}
\newtcolorbox{successbox}[1]{
  colback=UDsuccess!10,colframe=UDsuccess,
  fonttitle=\bfseries\small,title=#1,
  left=4pt,right=4pt,top=3pt,bottom=3pt
}

% ─── PROGRESO Y NUMERACIÓN ───────────────────────────────────
\setbeamertemplate{navigation symbols}{}
\setbeamertemplate{footline}{%
  \leavevmode
  \hbox{%
    \begin{beamercolorbox}[wd=.30\paperwidth,ht=2.5ex,dp=1.125ex,leftskip=2pt]{palette secondary}%
      \usebeamerfont{author in head/foot}\insertshortauthor
    \end{beamercolorbox}%
    \begin{beamercolorbox}[wd=.45\paperwidth,ht=2.5ex,dp=1.125ex,center]{palette tertiary}%
      \usebeamerfont{title in head/foot}\insertshorttitle
    \end{beamercolorbox}%
    \begin{beamercolorbox}[wd=.25\paperwidth,ht=2.5ex,dp=1.125ex,rightskip=4pt plus1fil]{palette primary}%
      \usebeamerfont{date in head/foot}\hfill
      \insertframenumber{} / \inserttotalframenumber
    \end{beamercolorbox}%
  }%
}

\setbeamertemplate{headline}{%
  \begin{beamercolorbox}[ht=1.6ex,dp=0.4ex,wd=\paperwidth]{palette quaternary}%
    \hspace{6pt}{\tiny\color{white}\insertsectionhead}
    \hfill
    \begin{tikzpicture}[overlay]
      \fill[UDaccent] (0,0) rectangle (\paperwidth,0.8pt);
    \end{tikzpicture}
  \end{beamercolorbox}%
}

% ─── BULLETS PERSONALIZADOS ──────────────────────────────────
\setbeamertemplate{itemize item}{\color{UDsecond}\tiny\raisebox{1pt}{$\blacksquare$}}
\setbeamertemplate{itemize subitem}{\color{UDprimary}\tiny$\triangleright$}
\setbeamertemplate{enumerate item}{\color{UDsecond}\bfseries\insertenumlabel.}

% ─── METADATOS ───────────────────────────────────────────────
\title[HPC: Potencial de Morse]{%
  \textbf{Optimización y Paralelización en HPC}\\
  \large Potencial de Morse — Cuatro Métodos Numéricos
}
\subtitle{Sistemas Distribuidos 2026-1 — Primer Parcial}
\author[Huertas · Rodriguez]{%
  \textbf{Camilo Huertas} \quad \textbf{Sebastián Rodriguez}\\
  {\small Programa Académico de Física}
}
\institute[UDFJC]{%
  Universidad Distrital Francisco José de Caldas\\
  {\small Docente: Carlos Andrés Gómez Vasco}
}
\date{27 de marzo de 2026}

% ============================================================
\begin{document}
% ============================================================

% ─── PORTADA ─────────────────────────────────────────────────
{
\setbeamertemplate{headline}{}
\setbeamertemplate{footline}{}
\begin{frame}
  \begin{tikzpicture}[remember picture, overlay]
    \fill[UDdark] (current page.south west) rectangle (current page.north east);
    \fill[UDprimary!80,opacity=0.7]
      ([yshift=-0.35\paperheight]current page.north west) rectangle
      ([yshift=0.35\paperheight]current page.south east);
    \fill[UDsecond]
      ([xshift=0pt]current page.south west) rectangle
      ([yshift=0.018\paperheight,xshift=\paperwidth]current page.south west);
    \node[anchor=north west,text=white,font=\tiny\itshape] at
      ([xshift=6pt,yshift=-2pt]current page.north west)
      {Universidad Distrital Francisco José de Caldas — Sistemas Distribuidos};
  \end{tikzpicture}
  \vspace{0.45cm}
  \titlepage
\end{frame}
}

% ─── TABLA DE CONTENIDO ──────────────────────────────────────
\begin{frame}{Agenda}
\begin{columns}[t]
\column{0.48\textwidth}
\tableofcontents[sections={1-6}]
\column{0.48\textwidth}
\tableofcontents[sections={7-12}]
\end{columns}
\end{frame}

% ============================================================
\section{Problema Físico}
% ============================================================

\begin{frame}{El Potencial de Morse: Descripción del Sistema}
\begin{columns}[c]
\column{0.50\textwidth}
\begin{physbox}{Modelo Molecular Diatómico}
El potencial de Morse describe la energía potencial de una molécula diatómica en función de la distancia internuclear $x$:
\[
  V(x) = D_e \left( 1 - e^{-a(x-x_e)} \right)^2
\]
\end{physbox}
\vspace{0.27cm}
\begin{block}{Parámetros Físicos Críticos}
\begin{itemize}\setlength\itemsep{2pt}
  \item $D_e$: Profundidad del pozo (energía de disociación)
  \item $x_e$: Distancia de equilibrio internuclear
  \item $a$: Parámetro de curvatura del enlace
  \item $\mu$: Masa reducida del sistema diatómico
\end{itemize}
\end{block}

\column{0.46\textwidth}
\begin{center}
\begin{tikzpicture}[scale=0.98]
  \begin{axis}[
    width=7.6cm, height=6.4cm,
    xlabel={$x$ (u.a.)},
    ylabel={$V(x)$},
    xmin=-1, xmax=5,
    ymin=-0.2, ymax=1.2,
    axis lines=left,
    grid=major, grid style={gray!20},
    tick label style={font=\tiny},
    label style={font=\small},
    legend style={font=\tiny, at={(0.95,0.5)}, anchor=east}
  ]
    % Morse potential
    \addplot[UDprimary, very thick, domain=-0.5:5, samples=100]
      {(1 - exp(-1.5*(x-1.0)))^2};
    \addplot[UDsecond!70, dashed, domain=-0.5:5, samples=100]
      {0.5*(x-1.0)^2};
    % Energy levels
    \addplot[UDaccent, thick, domain=0.1:3.5] {0.17};
    \addplot[UDaccent, thick, domain=0.0:4.0] {0.45};
    \addplot[UDaccent, thick, domain=-0.2:4.5] {0.67};
    \addplot[UDaccent, thick, domain=-0.4:4.8] {0.84};
    % Dissociation line
    \addplot[UDsecond, dashed, domain=-1:5] {1.0};
    \addlegendentry{Morse}
    \addlegendentry{Arm. Simple}
    \node[font=\tiny, UDsecond] at (axis cs:4.2,1.05) {$D_e$};
    \node[font=\tiny, UDaccent] at (axis cs:4.0,0.19) {$E_0$};
    \node[font=\tiny, UDaccent] at (axis cs:4.2,0.47) {$E_1$};
    \node[font=\tiny, UDaccent] at (axis cs:4.4,0.69) {$E_2$};
    \node[font=\tiny, UDaccent] at (axis cs:4.4,0.86) {$E_3$};
  \end{axis}
\end{tikzpicture}
\end{center}
\end{columns}
\end{frame}

% ─────────────────────────────────────────────────────────────
\begin{frame}[allowframebreaks]{Solución Analítica y Estados Acotados}
\begin{columns}[T]
\column{0.52\textwidth}
\begin{physbox}{Autovalores Exactos del Modelo de Morse}
\[
  E_n = \hbar\omega_e\!\left(n+\tfrac{1}{2}\right)
        - \frac{\left[\hbar\omega_e\!\left(n+\tfrac{1}{2}\right)\right]^2}{4D_e}
\]
donde $\omega_e = a\sqrt{2D_e/\mu}$ es la frecuencia de vibración clásica.
\end{physbox}
\vspace{-0.18cm}
\begin{alertbox}{\faExclamationTriangle\ Alcance del Estudio}
Se buscan \textbf{únicamente los primeros 4 niveles de energía} ($n=0,1,2,3$) para estandarizar la comparación entre todos los métodos numéricos.\\[2pt]
\textbf{Justificación:} Los estados superiores se acercan a $D_e$, donde divergen las aproximaciones de base finita y los métodos locales de baja resolución.
\end{alertbox}

\column{0.44\textwidth}
\begin{block}{Dos Propiedades Físicas Fundamentales}
\begin{enumerate}\setlength\itemsep{4pt}
  \item \textbf{Anarmonicidad:} Los niveles se comprimen con $n$ creciente
  \[
    \Delta E_{n} = \hbar\omega_e\left(1 - \frac{\hbar\omega_e(n+1)}{2D_e}\right)
  \]
  \item \textbf{Número finito de estados:} $n_{max} < \frac{D_e}{\hbar\omega_e} - \frac{1}{2}$
\end{enumerate}
\end{block}
\vspace{-0.155cm}
\begin{successbox}{Objetivo Computacional}
Calcular $\{E_0, E_1, E_2, E_3\}$ y cuantificar el \textbf{error absoluto} frente a la solución analítica exacta.
\end{successbox}
\end{columns}
\end{frame}

% ============================================================
\section{Migración Fortran → C}
% ============================================================

\begin{frame}[allowframebreaks]{Migración de Fortran 77 a C: Motivación y Objetivos}
\begin{columns}[T]
\column{0.50\textwidth}
\begin{alertbox}{\faBug\ Limitaciones del Código Fortran Original}
\begin{itemize}\setlength\itemsep{3pt}
  \item Subrutinas \texttt{DSYEV/DGEMM} sin interfaz moderna a LAPACKE
  \item Gestión de memoria estática (arreglos de tamaño fijo en \texttt{COMMON})
  \item Imposible garantizar alineación de datos para SIMD (AVX/AVX-512)
  \item Directivas OpenMP limitadas en Fortran 77 (\texttt{!OMP} rudimentario)
  \item Sin soporte para \texttt{restrict}, \texttt{aligned\_alloc} ni intrínsecos vectoriales
\end{itemize}
\end{alertbox}

\column{0.46\textwidth}
\begin{successbox}{\faRocket\ Ventajas del Ecosistema C + GCC}
\begin{itemize}\setlength\itemsep{3pt}
  \item \texttt{aligned\_alloc(64, ...)} $\rightarrow$ acceso AVX-512 sin penalidades
  \item \texttt{restrict}: el compilador vectoriza sin aliasing
  \item Interfaz directa a \texttt{LAPACKE\_dsyev} (C binding nativo)
  \item OpenMP 4.5+ con \texttt{collapse}, \texttt{guided}, \texttt{SIMD} clauses
  \item \texttt{-march=native -O3 -ffast-math} genera instrucciones \texttt{vmovapd}
\end{itemize}
\end{successbox}
\end{columns}

\vspace{0.27cm}
\begin{center}
\begin{tikzpicture}[
    node distance=1.2cm and 1.5cm, % vertical and horizontal distance
    auto, 
    font=\small,
    >=Stealth, % Tip de flecha global
    thick
]
  % --- PRIMERA FILA (Derecha) ---
  \node[draw,fill=UDsecond!20,rounded corners,minimum width=2.5cm,minimum height=0.7cm] (f77) {Fortran 77};
  
  \node[draw,fill=UDaccent!30,rounded corners,minimum width=2.5cm,minimum height=0.7cm, right=of f77] (anl) {Análisis HPC};
  
  \node[draw,fill=UDprimary!20,rounded corners,minimum width=2.5cm,minimum height=0.7cm, right=of anl] (cser) {C serial};

  % --- SEGUNDA FILA (Abajo y hacia la izquierda) ---
  % Bajamos desde cser
  \node[draw,fill=UDprimary!40,rounded corners,minimum width=2.5cm,minimum height=0.7cm, below=1.2cm of cser] (copt) {C optimizado};
  
  % Continuamos a la izquierda de copt
  \node[draw,fill=UDsuccess!40,rounded corners,minimum width=2.5cm,minimum height=0.7cm, left=of copt] (comp) {C + OpenMP};

  % --- ETIQUETAS (Relativas a los nodos) ---
  \node[above=0.05cm of f77, font=\tiny, UDsecond]  {Referencia};
  \node[above=0.05cm of anl, font=\tiny]           {Profiling};
  \node[above=0.05cm of cser, font=\tiny]          {Porteo};
  \node[below=0.05cm of copt, font=\tiny]          {Localidad};
  \node[below=0.05cm of comp, font=\tiny, UDsuccess] {Paralelismo};

  % --- CONEXIONES ---
  \draw[->, UDprimary] (f77) -- (anl);
  \draw[->, UDprimary] (anl) -- (cser);
  \draw[->, UDprimary] (cser) -- (copt); % Conexión vertical
  \draw[->, UDsuccess] (copt) -- (comp); % Conexión hacia la izquierda
\end{tikzpicture}
\end{center}
\end{frame}

% ============================================================
\section{Métodos Numéricos: Fundamento Físico}
% ============================================================

% --- DIAPOSITIVA 1: Teoría y Ecuaciones ---
\begin{frame}{Método 1 — Diagonalización en Base SHO (Teoría)}

\begin{physbox}{Ecuación de Schrödinger en Base Discreta}
Se resuelve $\hat{H}\ket{\psi} = E\ket{\psi}$ expandiendo en la base $\{\ket{n}\}$ del oscilador armónico lineal (LHO).
\end{physbox}

\vspace{0.4cm}

\textbf{Operadores en la Base de Fock:}

Los operadores $\hat{X}$ y $\hat{P}$ son tridiagonales:
\[
  X_{ij} , P_{ij} \propto  \frac{1}{\sqrt{2}}\!\left(\sqrt{j}\,\delta_{i,j-1} + \sqrt{i}\,\delta_{j,i-1}\right)
\]

\vspace{0.2cm}

\textbf{Construcción del Potencial:}
\[
  V_{ij} = \sum_{\gamma=1}^N (V_P)_{i\gamma}\, V(x_\gamma)\, (V_P)_{j\gamma}
\]

\vspace{0.2cm}

\textbf{Hamiltoniano Total:}
\[
  \mathbb{H} = \frac{1}{2\mu}\mathbb{P}^2 + \mathbb{V}
\]

\end{frame}


% --- DIAPOSITIVA 2: Estructura y Costo Computacional ---
% --- DIAPOSITIVA 2: Estructura y Costo Computacional ---
\begin{frame}{Método 1 — Diagonalización en Base SHO (Estructura y Costo)}

\begin{columns}[T] % Alineación superior

% --- Columna Izquierda: Gráfico TikZ ---
\column{0.5\textwidth}
\begin{center}
\begin{tikzpicture}[scale=0.85] 
  \draw[-Stealth,UDprimary,thick] (0,0) -- (4.2,0) node[right,font=\scriptsize]{$n$ (índice)};
  \draw[-Stealth,UDprimary,thick] (0,0) -- (0,4.2) node[above,font=\scriptsize]{$m$ (índice)};
  \foreach \i in {0,...,3} {
    \foreach \j in {0,...,3} {
      \pgfmathsetmacro{\val}{abs(\i-\j) < 2 ? 1 : 0}
      \ifnum\i=\j
        \fill[UDprimary!60] (\i*0.9+0.05,\j*0.9+0.05) rectangle (\i*0.9+0.85,\j*0.9+0.85);
      \else
        \ifnum\numexpr\i-\j\relax=1
          \fill[UDsecond!50] (\i*0.9+0.05,\j*0.9+0.05) rectangle (\i*0.9+0.85,\j*0.9+0.85);
        \else
          \ifnum\numexpr\j-\i\relax=1
            \fill[UDsecond!50] (\i*0.9+0.05,\j*0.9+0.05) rectangle (\i*0.9+0.85,\j*0.9+0.85);
          \else
            \fill[gray!10] (\i*0.9+0.05,\j*0.9+0.05) rectangle (\i*0.9+0.85,\j*0.9+0.85);
          \fi
        \fi
      \fi
    }
  }
  \node[font=\scriptsize,white] at (0.45,0.45) {$X_{00}$};
  \node[font=\scriptsize] at (1.85,2.35) {$\cdots$};
  \node[below,font=\scriptsize,UDprimary] at (2,-0.2) {Estructura Tridiagonal de $\hat{X}$};
  \draw[decorate,decoration={brace,amplitude=5pt},UDsecond,thick] (-0.1,3.7) -- (-0.1,-0.1) node[midway,left=3pt,font=\scriptsize,UDsecond]{$N$};
\end{tikzpicture}
\end{center}

% --- Columna Derecha: Caja de Costo Computacional ---
\column{0.5\textwidth}
\vspace{0.8cm} % Ajuste vertical para que quede centrada respecto al gráfico
\begin{alertbox}{Costo Computacional}
\begin{itemize}
  \item Construcción: $\mathcal{O}(N)$ (tridiagonal)
  \item Diagonalización: $\mathcal{O}(N^3)$ vía \texttt{DSYEV}
  \item Memoria: $\mathcal{O}(N^2)$
\end{itemize}
\end{alertbox}

\end{columns}

\end{frame}
% ─────────────────────────────────────────────────────────────
\begin{frame}{Método 2 — Diferencias Finitas (FDM): Fundamento}
\begin{columns}[T]
\column{0.50\textwidth}
\begin{physbox}{Discretización del Operador Cinético}
La derivada segunda se aproxima por diferencia central:
\[
  \left.\frac{d^2\psi}{dx^2}\right|_{x_i} \approx \frac{\psi_{i+1} - 2\psi_i + \psi_{i-1}}{\Delta x^2} + \mathcal{O}(\Delta x^2)
\]
\end{physbox}

\vspace{0.18cm}
\textbf{Hamiltoniano Matricial Tridiagonal:}
\[
  H_{ij} = \begin{cases}
    \dfrac{1}{\Delta x^2} + V(x_i) & i=j\\[6pt]
    -\dfrac{1}{2\Delta x^2} & i=j\pm 1\\[4pt]
    0 & \text{resto}
  \end{cases}
\]

\vspace{0.135cm}
\begin{block}{Convergencia y Estabilidad}
\begin{itemize}\setlength\itemsep{2pt}
  \item \textbf{Convergencia algebraica} $\mathcal{O}(\Delta x^2)$
  \item Requiere $N \gtrsim 1000$ para $10^{-5}$ de error relativo en $E_{0..3}$
  \item Costo: $\mathcal{O}(N^3)$ por diagonalización densa
\end{itemize}
\end{block}

\column{0.46\textwidth}
\begin{center}
\begin{tikzpicture}[scale=0.656]
  \begin{axis}[
    width=6cm, height=4cm,
    xlabel={$N$ (puntos de malla)},
    ylabel={Error relativo $|\Delta E/E|$},
    xmin=100, xmax=2000,
    ymin=1e-7, ymax=1e-1,
    ymode=log,
    xmode=log,
    axis lines=left,
    grid=major, grid style={gray!20},
    tick label style={font=\tiny},
    label style={font=\small},
    legend style={font=\tiny, at={(0.97,0.97)}, anchor=north east}
  ]
    \addplot[UDprimary, thick, mark=square*, mark size=2]
      coordinates {(100,5e-2)(200,1.3e-2)(500,2e-3)(1000,5e-4)(2000,1.2e-4)};
    \addplot[UDsecond, thick, mark=triangle*, mark size=2]
      coordinates {(100,5e-2)(200,1.3e-2)(500,2e-3)(1000,5e-4)(2000,1.2e-4)};
    \addplot[UDsuccess, thick, mark=*, mark size=2]
      coordinates {(100,5e-2)(200,1.3e-2)(500,2e-3)(1000,5e-4)(2000,1.2e-4)};
    \addplot[UDaccent!80!black, thick, mark=diamond*, mark size=2]
      coordinates {(100,6e-2)(200,1.5e-2)(500,2.5e-3)(1000,6e-4)(2000,1.5e-4)};
    \addlegendentry{$E_0$}
    \addlegendentry{$E_1$}
    \addlegendentry{$E_2$}
    \addlegendentry{$E_3$}
  \end{axis}
\end{tikzpicture}
\end{center}
\begin{alertbox}{Rango Óptimo FDM}
$\mathbf{N \in [1000, 2000]}$: error $< 10^{-4}$\\
Menor $N$: error de discretización $\uparrow$\\
Mayor $N$: costo $O(N^3)$ prohibitivo
\end{alertbox}
\end{columns}
\end{frame}

% ─────────────────────────────────────────────────────────────
% --- DIAPOSITIVA 1: Fundamento Físico y Numerov ---
\begin{frame}{Método 3 — Shooting + Numerov (Fundamento)}

\begin{physbox}{Ecuación de Schrödinger como IVP}
Se reescribe la ETID como EDO de 2do orden:
\[
  \frac{d^2\psi}{dx^2} = \hat{Q}(x,E)\psi,\quad \hat{Q} = -2\mu[E-V(x)]
\]
El autovalor $E$ es el \textbf{parámetro de control} que se ajusta iterativamente.
\end{physbox}

\vspace{0.4cm}

\textbf{Integrador de Numerov (error $\mathcal{O}(h^6)$) donde $k^2(x) = 2\mu(E-V(x))$.:}
\[
  \psi_{i+1} = \frac{2(1-\frac{5h^2}{12}k_i^2)\psi_i - (1+\frac{h^2}{12}k_{i-1}^2)\psi_{i-1}}{1+\frac{h^2}{12}k_{i+1}^2}
\]


\vspace{0.4cm}

\textbf{Teorema de Sturm-Liouville:}\\
El estado $n$-ésimo tiene \emph{exactamente} $n$ nodos. Esto permite una \textbf{búsqueda por bisección} guiada por el conteo de nodos en la función de onda.

\end{frame}


% --- DIAPOSITIVA 2: Visualización y Parámetros ---
\begin{frame}{Método 3 — Shooting + Numerov (Implementación)}

\begin{columns}[T]

% --- Columna Izquierda: Esquema del Shooting ---
\column{0.5\textwidth}
\begin{center}
\begin{tikzpicture}[scale=1] % Escala aumentada para mayor claridad
  % Energy axis
  \draw[-Stealth,UDprimary,thick] (0,-0.3) -- (0,4.5) node[above,font=\scriptsize]{$E$};
  % Potential schematic
  \fill[UDlight] (0.5,0) -- (0.5,0.8) -- (1.5,0.2) -- (2.5,0.2) -- (3.5,0.8) -- (3.5,0) -- cycle;
  \draw[UDprimary,thick] (0.3,3.5) parabola bend (2,0.1) (3.7,3.5);
  % Levels
  \foreach \lev/\col/\lab in {0.7/UDaccent/$E_0$, 1.4/UDaccent/$E_1$, 2.1/UDaccent/$E_2$, 2.7/UDaccent/$E_3$} {
    \draw[\col, thick] (0.8,\lev) -- (3.2,\lev);
    \node[right,font=\scriptsize,\col] at (3.2,\lev) {\lab};
  }
  % Shooting arrows
  \draw[-Stealth,UDsecond,thick,dashed] (-0.4,1.4) -- (0.8,1.4) node[left=12pt,font=\scriptsize,UDsecond]{Disparo};
  \node[font=\scriptsize] at (2.0,4.1) {¿$\psi(x_{max})\to 0$?};
  \draw[-Stealth,gray] (2.0,3.9) -- (2.0,3.0);
\end{tikzpicture}
\end{center}

% --- Columna Derecha: Caja de Rango y Complejidad ---
\column{0.45\textwidth}
\vspace{1.2cm}
\begin{alertbox}{Rango Óptimo Shooting}
\begin{itemize}
  \item $N_{steps} \in [500, 2000]$
  \item Tolerancia bisección: $10^{-10}$
  \item Complejidad: $\mathcal{O}(N_{steps} \times N_{iter})$
\end{itemize}
\end{alertbox}
\end{columns}

\end{frame}
% ─────────────────────────────────────────────────────────────
% --- DIAPOSITIVA 1: Fundamento y Matriz ---
\begin{frame}{Método 4 — Sinc-DVR (Fundamento y Operadores)}

% --- Bloque Superior (Ancho completo) ---
\begin{physbox}{Base de Funciones Cardinales Sinc}
La función de onda se expande en funciones de Whittaker (sinc):
\[
  C_j(x) = \text{sinc}\!\left(\frac{\pi(x-x_j)}{\Delta x}\right) = \frac{\sin[\pi(x-x_j)/\Delta x]}{\pi(x-x_j)/\Delta x}
\]
Ortonormalidad en puntos de malla: $\langle C_i|C_j\rangle = \Delta x\,\delta_{ij}$
\end{physbox}

\vspace{0.4cm}

% --- Bloque Inferior (Dividido en columnas/minipages) ---
\begin{columns}[T]

  % Columna Izquierda: Hamiltoniano
  \begin{column}{0.55\textwidth}
    \textbf{Hamiltoniano Sinc-DVR:}
    \small
    \[
      H_{ij} = \begin{cases}
        \dfrac{\pi^2}{6\Delta x^2} + V(x_i) & i=j\\[10pt]
        \dfrac{(-1)^{i-j}}{\Delta x^2(i-j)^2} & i\ne j
      \end{cases}
    \]
    \footnotesize
    Estructura \textbf{densa} y simétrica.
  \end{column}

  % Columna Derecha: Potencial
  \begin{column}{0.40\textwidth}
    \vspace{0.2cm}
    \textbf{Propiedad del Potencial:}
    \[ V_{ij}=V(x_i)\delta_{ij} \]
    
    \vspace{0.2cm}
    \footnotesize
    El potencial \textbf{permanece diagonal}, simplificando la construcción frente a otras bases espectrales que requieren integrales de solapamiento.
  \end{column}

\end{columns}

\end{frame}


% --- DIAPOSITIVA 2: Convergencia y Comparativa ---
\begin{frame}{Método 4 — Sinc-DVR (Eficiencia y Comparativa)}

\begin{columns}[T]

% --- Columna Izquierda: Éxito y Costo ---
\column{0.48\textwidth}
\begin{successbox}{Convergencia Espectral}
Caída \textbf{exponencial} del error con $N$ pequeño.\\
$N \sim 200$ $\Rightarrow$ error $< 10^{-6}$ para los primeros 4 estados ($E_{0..3}$).
\end{successbox}

\vspace{0.4cm}

\begin{alertbox}{Costo Computacional}
\begin{itemize}\setlength\itemsep{4pt}
  \item Construcción: $\mathcal{O}(N^2)$ (\textbf{densa})
  \item Diagonalización: $\mathcal{O}(N^3)$
  \item Memoria: $8N^2$ bytes
\end{itemize}
\end{alertbox}

% --- Columna Derecha: Tabla Comparativa ---
\column{0.48\textwidth}
\vspace{0.6cm}
\begin{block}{Comparativa vs. FDM}
\centering
\small
\begin{tabular}{lcc}
\toprule
\textbf{Característica} & \textbf{FDM} & \textbf{Sinc-DVR} \\
\midrule
Orden error & $\mathcal{O}(h^2)$ & Exponencial\\
Estructura $H$ & Rala & Densa\\
$N$ óptimo & 1000--2000 & 100--300\\
\bottomrule
\end{tabular}
\end{block}

\end{columns}

\end{frame}

% ─────────────────────────────────────────────────────────────
\begin{frame}{Rango Óptimo y Error de Cada Método (Primeros 4 Niveles)}
\begin{center}
\begin{tikzpicture}
  \begin{axis}[
    width=10.4cm, height=4.8cm,
    xlabel={Parámetro característico $N$},
    ylabel={Error relativo $|\Delta E / E_{analítico}|$},
    xmin=50, xmax=2100,
    ymin=1e-8, ymax=1e-0,
    ymode=log,
    axis lines=left,
    grid=major, grid style={gray!15},
    tick label style={font=\small},
    label style={font=\small},
    legend style={font=\small, at={(0.99,0.99)}, anchor=north east, fill=white, fill opacity=0.8},
    legend columns=2
  ]
    % Shooting
    \addplot[UDsecond,very thick,mark=square*,mark size=2.5,smooth]
      coordinates {(100,5e-3)(200,1e-3)(500,2e-5)(1000,2e-5)(2000,2e-5)};
    % FDM
    \addplot[UDprimary,very thick,mark=triangle*,mark size=2.5,smooth]
      coordinates {(100,8e-2)(200,2e-2)(500,3e-3)(1000,7e-4)(2000,1.5e-4)};
    % SHO Base
    \addplot[UDsuccess,very thick,mark=*,mark size=2.5,smooth]
      coordinates {(50,3e-4)(100,1e-5)(200,3e-6)(300,2e-6)(400,2e-6)};
    % Sinc-DVR
    \addplot[UDaccent!80!black,very thick,mark=diamond*,mark size=2.5,smooth]
      coordinates {(50,5e-3)(100,1e-5)(200,1e-7)(300,5e-8)(500,4e-8)};
    \addlegendentry{Shooting (Numerov)}
    \addlegendentry{Diferencias Finitas}
    \addlegendentry{Base SHO}
    \addlegendentry{Sinc-DVR}
    % Optimal regions
    \addplot[UDsecond!40, fill=UDsecond!10, fill opacity=0.5, draw=none]
      coordinates{(500,1e-8)(2000,1e-8)(2000,1e-0)(500,1e-0)} \closedcycle;
    \addplot[UDprimary!40, fill=UDprimary!10, fill opacity=0.5, draw=none]
      coordinates{(1000,1e-8)(2000,1e-8)(2000,1e-0)(1000,1e-0)} \closedcycle;
    \node[rotate=90,font=\tiny,UDsecond] at (axis cs:700,1e-4) {Óptimo Shoot};
    \node[rotate=90,font=\tiny,UDprimary] at (axis cs:1500,1e-3) {Óptimo FDM};
  \end{axis}
\end{tikzpicture}
\end{center}
\end{frame}

% ============================================================
\section{Implementaciones Seriales}
% ============================================================

\begin{frame}[fragile,allowframebreaks]{Serial: Base SHO en Fortran 77 (Código Original)}
\begin{columns}[T]
\column{0.48\textwidth}
\begin{block}{Construcción de Operadores $\hat{X}$, $\hat{P}$}
\begin{lstlisting}[style=fortranBeamer,language=Fortran]
C  BASE LHO: MATRICES X y P
   DO 10 I=1,N-1
      A_VAL=DBLE(I)
      X(I,I+1)=DSQRT(A_VAL/2.D0)
      X(I+1,I)=X(I,I+1)
      P(I,I+1)=-DSQRT(A_VAL/2.D0)
      P(I+1,I)=DSQRT(A_VAL/2.D0)
 10 CONTINUE
\end{lstlisting}
\end{block}
\vspace{-0.055cm}
\begin{block}{Potencial via Transformación Unitaria}
\begin{lstlisting}[style=fortranBeamer,language=Fortran]
C  V_ij = sum_k VP(i,k)*V(x_k)*VP(j,k)
   DO 20 I=1,N
    DO 30 J=I,N
       SUM=0.D0
       DO 40 K=1,N
         SUM=SUM+VP(I,K)*V_MORSE(X_EIG(K))*VP(J,K)
 40    CONTINUE
       VMAT(I,J)=SUM
       VMAT(J,I)=SUM
 30  CONTINUE
 20 CONTINUE
\end{lstlisting}
\end{block}

\column{0.48\textwidth}
\begin{block}{Diagonalización Final}
\begin{lstlisting}[style=fortranBeamer,language=Fortran]
C  H = T + V, extraer autovalores
   CALL MATMUT(P,P,T,N)
   DO 50 I=1,N
    DO 60 J=1,N
       H(I,J) = -0.5D0*T(I,J)+VMAT(I,J)
 60  CONTINUE
 50 CONTINUE
   CALL EIGEN(H,N,N,ENERGIA,VP,WORK1,WORK2,Z)
\end{lstlisting}
\end{block}
\vspace{-0.055cm}
\begin{alertbox}{Problema de Rendimiento}
\begin{itemize}
  \item Triple bucle $O(N^3)$ para construir $V_{ij}$
  \item \texttt{MATMUT} calcula $P^2$ como producto denso
  \item Acceso con \textbf{stride-$N$}: falla sistemática de caché L1
\end{itemize}
\end{alertbox}
\end{columns}
\end{frame}

% ─────────────────────────────────────────────────────────────
\begin{frame}[fragile,allowframebreaks]{Serial C: Base SHO — Explotación de Raleza}
\begin{columns}[T]
\column{0.50\textwidth}
\begin{block}{Operador Cinético Pentadiagonal $O(N)$}
\begin{lstlisting}[language=C,style=cppBeamer]
// Solo se computan los 5 diagonales de T = P^2
for (int i = 0; i < n; i++) {
  if (i > 0)
    T[i*n+i] += P[i*n+(i-1)]*P[(i-1)*n+i];
  if (i < n-1)
    T[i*n+i] += P[i*n+(i+1)]*P[(i+1)*n+i];
  if (i < n-2) {
    double val = P[i*n+(i+1)]*P[(i+1)*n+(i+2)];
    T[i*n+(i+2)] = T[(i+2)*n+i] = val;
  }
}
\end{lstlisting}
\end{block}

\begin{successbox}{Impacto}
$\mathcal{O}(N^3) \rightarrow \mathcal{O}(N)$\\
Speedup serial: $\approx\mathbf{2.5\times}$
\end{successbox}

\column{0.46\textwidth}
\begin{block}{Transformación de Base con Localidad Stride-1}
\begin{lstlisting}[language=C,style=cppBeamer]
// W = VP * diag(V(E)) — pre-computo
for(int i=0; i<n; i++)
  for(int k=0; k<n; k++)
    W[i*n+k] = VP[i*n+k] * VE[k];
// VMAT con acceso contiguo (Row-Major)
for (int i=0; i<n; i++) {
  for (int j=i; j<n; j++) {
    double sum = 0.0;
    const double *ri = &W[i*n];
    const double *rj = &VP[j*n];
    for (int k=0; k<n; k++)
      sum += ri[k] * rj[k];

\end{lstlisting}
\end{block}
\begin{block}{Técnica Clave}
Matriz intermedia $W_{ik}$ $\Rightarrow$ bucle interno con \textbf{stride-1}, habilitando vectorización AVX automática.
\end{block}
\end{columns}
\end{frame}

% ─────────────────────────────────────────────────────────────
\begin{frame}[fragile,allowframebreaks]{Serial: FDM en Fortran 77 y Versión C Optimizada}
\begin{columns}[T]
\column{0.48\textwidth}
\begin{block}{Fortran 77 — Código Original}
\begin{lstlisting}[style=fortranBeamer,language=Fortran]
C  Hamiltoniano FDM tridiagonal
   DX2 = DX * DX
   DO 100 I=1,N
      XI = XMIN + DBLE(I)*DX
      VI = D_POT*(1.D0-DEXP(-ALPHA*(XI-RE)))**2
      H(I,I) = (1.D0/DX2) + VI
      IF (I .LT. N) THEN
         H(I,I+1) = -1.D0/(2.D0*DX2)
         H(I+1,I) = H(I,I+1)
      END IF
100 CONTINUE
\end{lstlisting}
\end{block}
\begin{alertbox}{Problema}
\begin{itemize}
  \item \texttt{DEXP} evaluado $N$ veces con división interna
  \item Almacena en matriz densa aunque es tridiagonal, sin precomputo
  \item Sin pre-cómputo de constantes
\end{itemize}
\end{alertbox}

\column{0.48\textwidth}
\begin{block}{C Optimizado — Eliminación de Redundancias}
\begin{lstlisting}[language=C,style=cppBeamer]
const double dx2_inv = 1.0/(dx*dx);
const double off_diag = -0.5*dx2_inv;
for (int i=0; i<n; i++) {
  const double xi = XMIN + (double)(i+1)*dx;// exp(-alpha*(xi-RE)) — RE absorbido en XMIN
  const double et = exp(-ALPHA*xi);
  const double br = 1.0 - et;
  const double vi = D_POT*(br*br);
  H[i*n+i] = dx2_inv + vi;
  if (i < n-1) {
    H[i*n+(i+1)] = off_diag;
    H[(i+1)*n+i] = off_diag;
 
\end{lstlisting}
\end{block}
\vspace{-0.155cm}
\begin{successbox}{Speedup serial FDM}
\textbf{$> 100\times$ de mejora asintótica}\\
Factor principal: eliminación de potencias redundantes y uso de constantes pre-calculadas.
\end{successbox}
\end{columns}
\end{frame}

% ─────────────────────────────────────────────────────────────
\begin{frame}[fragile,allowframebreaks]{Serial: Shooting — Fortran 77 y Memoización en C}
\begin{columns}[T]
\column{0.48\textwidth}
\begin{block}{Fortran 77 — Bisección y Control de Nodos}
\begin{lstlisting}[style=fortranBeamer,language=Fortran]
C  Biseccion con teorema de Sturm
   DO WHILE(DABS(EUP-ELOW).GT.TOL)
      ETRIAL=(EUP+ELOW)/2.D0
      CALL PROPAGAR(ETRIAL,PSI_FINAL,NODOS)

      IF(NODOS.GT.N_TARGET) THEN
         EUP=ETRIAL
      ELSE IF(NODOS.LT.N_TARGET) THEN
         ELOW=ETRIAL
      ELSE
         IF(PSI_FINAL*SIGNO_REF.GT.0.D0) THEN
            EUP=ETRIAL
         ELSE
            ELOW=ETRIAL
         END IF
      END IF
   END DO
\end{lstlisting}
\end{block}

\column{0.48\textwidth}
\begin{block}{C Optimizado — Tabulación del Potencial}
\begin{lstlisting}[language=C,style=cppBeamer]
// Memoizacion: V(x) se calcula UNA sola vez
for (int i=0; i<=nsteps; i++) {
  double x = xmin + i*dx;
  double tmp = 1.0 - exp(-alpha*x);
  // 2*V para eliminar multiplicacion en nucleo
  v_base[i] = 2.0*d_pot*tmp*tmp;
  if (i < nsteps) {
    double xm = x + 0.5*dx;
    double tm = 1.0-exp(-alpha*xm);
    v_half[i] = 2.0*d_pot*tm*tm;
  }
}
// Nucleo RK4 accede stride-1 a v_base/v_half
for (int i=0; i<n; i++) {
  const double rk1_phi = (v_b[i]-en2)*psi;
  const double vm_term = v_h[i]-en2;
  const double rk2_phi = vm_term*(psi+h_2*rk1_psi);
  const double rk3_phi = vm_term*(psi+h_2*rk2_psi);
  const double rk4_phi = (v_b[i+1]-en2)*(psi+h*rk3_psi);
  psi += h_6*(rk1_psi+2*rk2_psi+2*rk3_psi+rk4_psi);
}
\end{lstlisting}
\end{block}
\begin{successbox}{Speedup: $\mathbf{3.5\times}$ serial}
\end{successbox}
\end{columns}
\end{frame}

% ─────────────────────────────────────────────────────────────
\begin{frame}[fragile,allowframebreaks]{Serial: Sinc-DVR — Fortran 77 y Optimización de Micro-Arquitectura}
\begin{columns}[T]
\column{0.48\textwidth}
\begin{block}{Fortran 77 — Construcción Original}
\begin{lstlisting}[style=fortranBeamer,language=Fortran]
C  Hamiltoniano Sinc-DVR denso
   DO 200 I=1,N
      XI=XMIN+DBLE(I)*DX
      VI=D_POT*(1.D0-DEXP(-ALPHA*(XI-RE)))**2
      DO 300 J=1,N
         IF (I.EQ.J) THEN
            H(I,J)=(PI**2)/(6.D0*DX**2)+VI
         ELSE
            SIGN_FACTOR=(-1.D0)**(I-J)
            DIST2=DBLE(I-J)**2
            H(I,J)=SIGN_FACTOR/(DX**2*DIST2)
         END IF
300   CONTINUE
200 CONTINUE
\end{lstlisting}
\end{block}
\begin{alertbox}{Problema de rendimiento}
\texttt{(-1.D0)**(I-J)} y \texttt{**2} son llamadas a \texttt{pow()}: \textbf{40 ciclos} vs. 1 ciclo con bit AND.
\end{alertbox}

\column{0.48\textwidth}
\begin{block}{C Optimizado — Bit-Trick y Triángulo Superior}
\begin{lstlisting}[language=C,style=cppBeamer]
const double diag_kin = (M_PI*M_PI)/(6.0*dx2);
for (int i=0; i<n; i++) {
  double xi = XMIN + (double)(i+1)*dx;
  double vi = D_POT*pow(1.0-exp(-ALPHA*xi),2);
  H[i*n+i] = diag_kin + vi;
  for (int j=i+1; j<n; j++) {  // Solo triangulo superior => 50% menos escrituras
    const int diff = i-j;
    const double inv_d2 = 1.0/(double)(diff*diff);
    // Bit AND: paridad en 1 ciclo vs. 40 ciclos de pow()
    const double sign = ((i-j)&1) ? -1.0 : 1.0;
    H[i*n+j] = (sign*dx2_inv)*inv_d2;
\end{lstlisting}
\end{block}
\begin{block}{Técnicas Aplicadas}
\begin{itemize}\setlength\itemsep{2pt}
  \item Eliminación de \texttt{if} dentro del bucle interno
  \item \texttt{aligned\_alloc(64,...)} $\Rightarrow$ instrucciones \texttt{vmovapd}
  \item Speedup: $\mathbf{1.8\times - 2.2\times}$
\end{itemize}
\end{block}
\end{columns}
\end{frame}

% ============================================================
\section{Banderas de Compilación}
% ============================================================

\begin{frame}[allowframebreaks]{Análisis de Banderas: Metodología y Métricas}
\begin{columns}[T]
\column{0.50\textwidth}
\begin{physbox}{Protocolo Experimental}
Se evaluaron 4 niveles de optimización en GCC:
\begin{enumerate}
  \item \textbf{-O0}: Sin optimización (código de referencia)
  \item \textbf{-O1}: Eliminación de redundancias básicas
  \item \textbf{-O2}: Vectorización SSE (registros XMM 128-bit)
  \item \textbf{-O3}: Vectorización agresiva AVX (registros YMM 256-bit) + \texttt{-march=native}
\end{enumerate}
\end{physbox}
\vspace{-0.18cm}
\begin{block}{Tres Métricas Capturadas}
\begin{enumerate}\setlength\itemsep{3pt}
  \item \textbf{Tamaño del segmento \texttt{.text}} 
  \item \textbf{Conteo de saltos} (branches): 
  \item \textbf{Reportes de inlining y vectorización} (\texttt{-fopt-info-vec})
\end{enumerate}
\end{block}

\column{0.46\textwidth}
\begin{block}{Resultados: Saltos y Tamaño del Binario}
\begin{center}
\small
\begin{tabular}{lrrrr}
\toprule
\textbf{Método} & \textbf{Saltos} & \textbf{Saltos} & \textbf{Bytes} & \textbf{Bytes}\\
 & \textbf{-O0} & \textbf{-O3} & \textbf{-O0} & \textbf{-O3}\\
\midrule
\texttt{morse\_base}  & 52 & 73 & 4636 & 4729\\
\texttt{morse\_fdm}   & 23 & 27 & 3771 & 3627\\
\texttt{morse\_sinc}  & 27 & 27 & 3626 & 3245\\
\texttt{morse\_shoot} & 24 & 27 & 4300 & 4328\\
\bottomrule
\end{tabular}
\end{center}
\end{block}
\vspace{-0.135cm}
\begin{alertbox}{Interpretación}
El aumento de saltos en \texttt{morse\_base} (52$\to$73) indica \textbf{loop versioning}: el compilador genera rutas alternativas del bucle para manejar restos de vectorización SIMD.
\end{alertbox}
\end{columns}
\end{frame}

% ─────────────────────────────────────────────────────────────
% --- DIAPOSITIVA 1: Visualización de Tiempos ---
\begin{frame}{Resultados de Banderas: Comparativa Visual}

\begin{columns}[T]

% --- Columna Izquierda: Gráfica de Barras ---
\column{0.60\textwidth}
\begin{center}
\begin{tikzpicture}
  \begin{axis}[
    ybar, bar width=10pt, % Aumenté un poco el ancho de barra
    width=7.5cm, height=5.5cm, % Gráfica más grande y legible
    symbolic x coords={-O0,-O1,-O2,-O3},
    xtick=data,
    ylabel={Tiempo relativo (base -O0)},
    ymin=0, ymax=1.15,
    axis lines=left,
    grid=major, grid style={gray!15},
    tick label style={font=\small},
    label style={font=\small},
    legend style={font=\tiny, at={(0.98,0.98)}, anchor=north east},
    legend columns=2
  ]
    \addplot[fill=UDsecond!70, draw=UDsecond]
      coordinates {(-O0,1.0)(-O1,0.38)(-O2,0.29)(-O3,0.25)};
    \addplot[fill=UDprimary!70, draw=UDprimary]
      coordinates {(-O0,1.0)(-O1,0.42)(-O2,0.35)(-O3,0.34)};
    \addplot[fill=UDsuccess!70, draw=UDsuccess]
      coordinates {(-O0,1.0)(-O1,0.55)(-O2,0.50)(-O3,0.48)};
    \addplot[fill=UDaccent!70, draw=UDaccent!70!black]
      coordinates {(-O0,1.0)(-O1,0.65)(-O2,0.61)(-O3,0.59)};
    \addlegendentry{Shooting}
    \addlegendentry{FDM}
    \addlegendentry{Base SHO}
    \addlegendentry{Sinc-DVR}
  \end{axis}
\end{tikzpicture}
\end{center}

% --- Columna Derecha: El Salto Principal ---
\column{0.35\textwidth}
\vspace{1.5cm}
\begin{successbox}{Mayor salto: -O0 $\to$ -O1}
La ganancia más drástica ocurre al activar las optimizaciones básicas:
\begin{itemize}
    \item Eliminación de código muerto.
    \item Reducción de llamadas redundantes al stack.
\end{itemize}
\end{successbox}

\end{columns}
\end{frame}


% --- DIAPOSITIVA 2: Análisis Técnico y Cuellos de Botella ---
\begin{frame}{Resultados de Banderas: Análisis y Límites}

\begin{columns}[T]

% --- Columna Izquierda: Análisis por Método ---
\column{0.58\textwidth}
\begin{block}{Análisis por Método}
\small
\begin{itemize}\setlength\itemsep{4pt}
  \item \textbf{Shooting} ($\Delta_{\max}$): Mayor beneficio en \texttt{-O3} por su naturaleza aritmética escalar e \emph{inlining}.
  \item \textbf{FDM}: \emph{Loop split} mejora la predicción de saltos en la malla densa.
  \item \textbf{Sinc-DVR}: Converge en \texttt{-O2}. Estructura ya compacta; \texttt{-O3} no aporta redundancias adicionales.
  \item \textbf{Base SHO}: En \texttt{-O3} aparece \emph{loop versioning} para registros XMM en la transformación.
\end{itemize}
\end{block}

% --- Columna Derecha: Alerta de Memoria ---
\column{0.38\textwidth}
\vspace{1.0cm}
\begin{alertbox}{Límite Detectado}
\textbf{Memory-bound}:\\
\vspace{0.2cm}
\footnotesize
FDM y Sinc-DVR alcanzan un cuello de botella en el ancho de banda de memoria desde \texttt{-O2}. El procesador queda en espera de datos de la RAM.
\end{alertbox}

\end{columns}
\end{frame}

% ─────────────────────────────────────────────────────────────
\begin{frame}[allowframebreaks]{Logs de Compilación: Inlining, Vectorización y Loop Splitting}
\begin{columns}[T]
\column{0.50\textwidth}
\begin{block}{\faSearch\ Fenómenos Observados en los Logs}
\begin{enumerate}\setlength\itemsep{4pt}
  \item \textbf{Inlining de \texttt{v\_morse}:}\\
  En \texttt{morse\_shoot}, desde \texttt{-O1} el compilador inserta el cuerpo de la función directamente en el integrador. Elimina el overhead en RK4.

  \item \textbf{Vectorización de bucles diagonales:}\\
  Log reporta \texttt{"loop vectorized using 16/32 byte vectors"} en \texttt{-O2/-O3}. : múltiples elementos de la diagonal calculados en 1 ciclo.

  \item \textbf{Loop Splitting en FDM:}\\
  Bajo \texttt{-O3}, el compilador divide el bucle de construcción del Hamiltoniano en 2 fragmentos: uno para la diagonal y otro para los elementos off-diagonal. Mejora la predicción de saltos y la localidad L1.
\end{enumerate}
\end{block}

\column{0.46\textwidth}
\begin{alertbox}{Tiling (Blocking) para Caché}
Se implementó \textit{tiling} en la transformación de base del Hamiltoniano:
\begin{itemize}
  \item FFT: tiempo reducido de $1.94\,$s $\to$ $0.64\,$s
  \item Métodos locales (Base/Lanczos): impacto mínimo; el conjunto de datos $N=400$ ya cabe en la caché L1D (2 MiB del AMD Ryzen Threadripper)
\end{itemize}
\end{alertbox}
\vspace{0.025cm}
\begin{conclusionbox}{Conclusión del Análisis de Compilación}
La optimización \textbf{no es solo una bandera de compilador}
\end{conclusionbox}
\end{columns}
\end{frame}

% ============================================================
\section{Serial vs. Serial Optimizado}
% ============================================================

% --- DIAPOSITIVA 1: Tabla Comparativa de Rendimiento ---
\begin{frame}{Comparativa: Serial Original vs. Serial Optimizado}

\begin{center}
\small % Reducimos ligeramente el tamaño para que la tabla respire mejor
\begin{tabular}{lp{3.2cm}p{3.2cm}cc}
\toprule
\textbf{Método} & \textbf{Optimización Principal} & \textbf{Factor Limitante} & \textbf{Speedup} & \textbf{$T_{serial}$ (s)}\\
\midrule
\textbf{FDM} & Pre-cómputo; eliminar \texttt{pow()} & Ancho de banda L3 & $> 100\times$ & 7.191 \\
\textbf{Shooting} & Memoización; núcleo RK4 stride-1 & Evaluación exponencial & $3.5\times$ & 0.232 \\
\textbf{Base SHO} & Op. $\hat{T}$ pentadiagonal $O(N)$ & Latencia (stride-N) & $2.5\times$ & 0.179 \\
\textbf{Sinc-DVR} & Bit-AND; triángulo superior & Diag. densa LAPACK & $1.8$--$2.2\times$ & 0.110 \\
\bottomrule
\end{tabular}
\end{center}

\vspace{0.4cm}
\footnotesize
\textbf{Nota:} Los tiempos corresponden a la ejecución serial optimizada comparada con la versión base sin banderas de optimización ni refinamiento de algoritmos.

\end{frame}


% --- DIAPOSITIVA 2: Análisis de Casos y Cuellos de Botella ---
\begin{frame}{Análisis de Optimización: Éxitos y Limitaciones}

\begin{columns}[T]

% Columna 1: El éxito de FDM
\column{0.48\textwidth}
\begin{successbox}{FDM: Caso Extremo}
\small
La mejora de $100\times$ se debe a la eliminación de cálculos redundantes de \texttt{DEXP} y \texttt{**2} dentro del bucle crítico. Al pre-computar constantes, el Hamiltoniano esparso se explota eficientemente, pasando de una carga pesada de CPU a una limitada por memoria.
\end{successbox}

% Columna 2: El límite de Sinc-DVR
\column{0.48\textwidth}
\begin{alertbox}{Sinc-DVR: Límite Algorítmico}
\small
Aunque la construcción de $H_{ij}$ es eficiente ($O(N^2)$), el tiempo total está dominado por la diagonalización densa vía \texttt{LAPACKE\_dsyev} ($O(N^3)$). Aquí la optimización de código tiene un techo; el cuello de botella es el algoritmo de diagonalización mismo.
\end{alertbox}

\end{columns}

\vspace{0.5cm}
\begin{block}{Conclusión Técnica}
\footnotesize
La transición de un código "orientado a cálculo" a uno "orientado a memoria" es evidente en FDM, mientras que en métodos espectrales como Sinc-DVR, la complejidad algorítmica prevalece sobre la implementación.
\end{block}

\end{frame}
% ─────────────────────────────────────────────────────────────
\begin{frame}{Gráficas de Comparativa Serial: Cuatro Métodos}
\begin{columns}[T]
\column{0.50\textwidth}
\begin{center}
\begin{tikzpicture}
  \begin{axis}[
    width=6cm, height=4cm,
    title={\small Comparativa Serial vs Optimizado},
    xlabel={$N$}, ylabel={Tiempo (s)},
    xmode=log, ymode=log,
    xmin=50, xmax=2000,
    ymin=1e-4, ymax=20,
    grid=major, grid style={gray!15},
    tick label style={font=\tiny},
    label style={font=\small},
    legend style={font=\tiny,at={(0.97,0.03)},anchor=south east},
    legend columns=2
  ]
    % FDM serial
    \addplot[UDprimary, thick, dashed, mark=triangle,mark size=2]
      coordinates{(100,0.03)(200,0.25)(500,3.6)(1000,7.2)(2000,15)};
    % FDM opt
    \addplot[UDprimary, very thick, mark=triangle*,mark size=2]
      coordinates{(100,3e-4)(200,2e-3)(500,2.5e-2)(1000,0.04)(2000,0.09)};
    % Shooting serial
    \addplot[UDsecond, thick, dashed, mark=square,mark size=2]
      coordinates{(500,0.01)(1000,0.07)(2000,0.23)};
    % Shooting opt
    \addplot[UDsecond, very thick, mark=square*,mark size=2]
      coordinates{(500,4e-3)(1000,2e-2)(2000,0.07)};
    \addlegendentry{FDM Orig.}
    \addlegendentry{FDM Opt.}
    \addlegendentry{Shoot Orig.}
    \addlegendentry{Shoot Opt.}
  \end{axis}
\end{tikzpicture}
\end{center}

\column{0.46\textwidth}
\begin{center}
\begin{tikzpicture}
  \begin{axis}[
    width=5.6cm, height=4cm,
    title={\small Base SHO y Sinc-DVR},
    xlabel={$N$}, ylabel={Tiempo (s)},
    xmode=log, ymode=log,
    xmin=50, xmax=500,
    ymin=1e-4, ymax=1,
    grid=major, grid style={gray!15},
    tick label style={font=\tiny},
    label style={font=\small},
    legend style={font=\tiny,at={(0.97,0.03)},anchor=south east},
    legend columns=2
  ]
    % Base serial
    \addplot[UDsuccess, thick, dashed, mark=o,mark size=2]
      coordinates{(50,3e-3)(100,1.2e-2)(200,0.09)(400,0.18)};
    % Base opt
    \addplot[UDsuccess, very thick, mark=*,mark size=2]
      coordinates{(50,1.5e-3)(100,5e-3)(200,0.04)(400,0.11)};
    % Sinc serial
    \addplot[UDaccent!80!black, thick, dashed, mark=diamond,mark size=2]
      coordinates{(50,3e-3)(100,1.1e-2)(200,0.06)(500,0.11)};
    % Sinc opt
    \addplot[UDaccent!80!black, very thick, mark=diamond*,mark size=2]
      coordinates{(50,1.8e-3)(100,6e-3)(200,0.03)(500,0.11)};
    \addlegendentry{SHO Orig.}
    \addlegendentry{SHO Opt.}
    \addlegendentry{DVR Orig.}
    \addlegendentry{DVR Opt.}
  \end{axis}
\end{tikzpicture}
\end{center}
\end{columns}

\vspace{0.135cm}
\begin{conclusionbox}{Síntesis: El Cuello de Bottleneck no es el Hardware}
Las mejoras más dramáticas vienen de \textbf{decisiones algorítmicas}: explotación de raleza, tabulación del potencial y acceso contiguo a memoria, no de banderas del compilador.
\end{conclusionbox}
\end{frame}

% ============================================================
\section{Paralelización con OpenMP}
% ============================================================

\begin{frame}[allowframebreaks]{¿Por qué OpenMP? Justificación desde la Arquitectura}
\begin{columns}[T]
\column{0.50\textwidth}
\begin{physbox}{AMD Ryzen Threadripper — Arquitectura Target}
\begin{itemize}\setlength\itemsep{3pt}
  \item 8 núcleos físicos con hiperthreading
  \item Caché L1D: 2 MiB por núcleo (datos)
  \item Caché L2: 512 KB por núcleo

  \item for-chanel  ancho de banda $\sim$100 GB/s
\end{itemize}
\end{physbox}
\vspace{-0.18cm}
\begin{block}{Diagnóstico Post-Optimización Serial}
A pesar de la optimización, los métodos matriciales (FDM y Sinc-DVR) alcanzan un \textbf{techo de memoria (memory-bound)}: el procesador está inactivo esperando datos de RAM. La solución: \textbf{dividir el dominio} para que cada núcleo opere sobre un subconjunto que cabe en su caché L1/L2.
\end{block}

\column{0.46\textwidth}
\begin{successbox}{Condiciones para OpenMP}
\begin{enumerate}\setlength\itemsep{4pt}
  \item \textbf{Paralelismo de datos}: bucles independientes sobre filas/columnas de matrices
  \item \textbf{Paralelismo de tareas}: búsqueda independiente de autovalores (Shooting)
  \item \textbf{Sin condiciones de carrera}: cada hilo accede a regiones de memoria exclusivas (excepto en sección crítica de Shooting)
\end{enumerate}
\end{successbox}
\vspace{0.18cm}

\end{columns}
\end{frame}

% ============================================================
\section{Implementaciones OpenMP}
% ============================================================

% --- DIAPOSITIVA 1: Afinidad NUMA y Operador Cinético ---
\begin{frame}[fragile]{OMP: Base SHO — Política NUMA y Operador Cinético}

\begin{columns}[T]

% Columna 1: First-Touch
\column{0.48\textwidth}
\begin{block}{First-Touch Policy (Afinidad NUMA)}
\begin{lstlisting}[language=C,style=cppBeamer,basicstyle=\tiny\ttfamily]
// Cada hilo inicializa su propia region
// => paginas fisicas en el nodo local
#pragma omp parallel for schedule(static)
for(int i=0; i<n2; i++) {
  X[i]=0.0; P[i]=0.0;
  T[i]=0.0; VMAT[i]=0.0;
}
\end{lstlisting}
\end{block}
\footnotesize
Garantiza que la memoria esté cerca del núcleo que la procesa, reduciendo la latencia de acceso en sistemas multiprocesador.

% Columna 2: Operador Cinético
\column{0.48\textwidth}
\begin{block}{Operador Cinético (Static)}
\begin{lstlisting}[language=C,style=cppBeamer,basicstyle=\tiny\ttfamily]
// Carga uniforme => static scheduling
#pragma omp parallel for schedule(static)
for (int i=0; i<n; i++) {
  if (i>0)
    T[i*n+i] += P[i*n+(i-1)]*P[(i-1)*n+i];
  if (i<n-1)
    T[i*n+i] += P[i*n+(i+1)]*P[(i+1)*n+i];
  if (i<n-2) {
    double val=P[i*n+(i+1)]*P[(i+1)*n+(i+2)];
    T[i*n+(i+2)] = T[(i+2)*n+i] = val;
  }
}
\end{lstlisting}
\end{block}
\end{columns}

\end{frame}


% --- DIAPOSITIVA 2: Transformación de Base y Scheduling ---
\begin{frame}[fragile]{OMP: Base SHO — Transformación y Balance de Carga}

\begin{columns}[T]

% Columna 1: Código de Transformación
\column{0.62\textwidth}
\begin{block}{Transformación con Guided Scheduling}
\begin{lstlisting}[language=C,style=cppBeamer,basicstyle=\tiny\ttfamily]
// W = VP*diag(V): collapse para 2 bucles
#pragma omp for schedule(static) collapse(2)
for(int i=0; i<n; i++)
  for(int k=0; k<n; k++)
    W[i*n+k] = VP[i*n+k]*VE[k];

// VMAT triangular: guided para desbalance
#pragma omp for schedule(guided)
for (int i=0; i<n; i++) {
  for (int j=i; j<n; j++) {
    double sum=0.0;
    const double *ri = &W[i*n];
    const double *rj = &VP[j*n];
    #pragma GCC ivdep
    for (int k=0; k<n; k++) sum+=ri[k]*rj[k];
    VMAT[i*n+j] = VMAT[j*n+i] = sum;
  }
}
\end{lstlisting}
\end{block}

% Columna 2: Explicación del Scheduling
\column{0.34\textwidth}
\vspace{0.5cm}
\begin{alertbox}{\texttt{collapse} + \texttt{guided}}
\small
\textbf{Collapse(2):} Aplana el espacio de iteración para mejorar la granularidad.

\vspace{0.2cm}
\textbf{Guided:} Compensa el desbalance natural al calcular solo el triángulo superior de la matriz simétrica.
\end{alertbox}

\end{columns}

\end{frame}

% ─────────────────────────────────────────────────────────────
\begin{frame}[fragile]{OMP: FDM — Static Scheduling y Aislamiento BLAS}
\begin{columns}[T]
\column{0.50\textwidth}
\begin{block}{Inicialización Paralela NUMA}
\begin{lstlisting}[language=C,style=cppBeamer]
// First-Touch: distribuye la memoria entre nodos
#pragma omp parallel for schedule(static)
for (size_t i=0; i<n2; i++)
  H[i] = 0.0;
\end{lstlisting}
\end{block}
\vspace{0.18cm}
\begin{block}{Construcción Paralela del Hamiltoniano}
\begin{lstlisting}[language=C,style=cppBeamer]
// Carga uniforme => static scheduling optimo
#pragma omp parallel for schedule(static)
for (int i=0; i<n; i++) {
  const double xi = XMIN+(double)(i+1)*dx;
  const double et = exp(-ALPHA*xi);
  const double vi = D_POT*(1.0-et)*(1.0-et);

  H[i*n+i] = dx2_inv + vi;
  if (i<n-1) {
    H[i*n+(i+1)] = off_diag;
    H[(i+1)*n+i] = off_diag;
  }
}
\end{lstlisting}
\end{block}

\column{0.46\textwidth}
\begin{block}{Aislamiento de BLAS/LAPACK}
\begin{lstlisting}[language=bash,style=cppBeamer]
# Evitar sobresubscripcion de nucleos
export OPENBLAS_NUM_THREADS=1
./morse_fdm_omp <N_puntos> <N_hilos>
\end{lstlisting}
\end{block}
\vspace{0.18cm}
\begin{alertbox}{Límite de FDM con OpenMP}
La diagonalización \texttt{LAPACKE\_dsyev} es $O(N^3)$ y se ejecuta de forma serializada. El Speedup de OpenMP se limita a la fase de \textbf{construcción} del Hamiltoniano.\\[4pt]
Para $N=2000$: $T_{constr} \ll T_{diag}$\\
$\Rightarrow$ Speedup global $S \approx 1.02$
\end{alertbox}
\vspace{0.135cm}
\begin{conclusionbox}{Sin Race Conditions}
Cada hilo opera sobre fila $i$ exclusiva. No se necesitan secciones críticas ni reducción.
\end{conclusionbox}
\end{columns}
\end{frame}

% ─────────────────────────────────────────────────────────────
\begin{frame}[fragile]{OMP: Shooting — Paralelismo de Tareas y Sincronización}

% --- Parte Superior: Dos Columnas de Código ---
\begin{columns}[T]

  % Columna Izquierda: Búsqueda Concurrente
  \begin{column}{0.48\textwidth}
    \begin{block}{Búsqueda Concurrente}
    \begin{lstlisting}[language=C,style=cppBeamer,basicstyle=\tiny\ttfamily]
// Dynamic: biseccion heterogenea
#pragma omp parallel for schedule(dynamic)
for (int i=0; i<n_intervals; i++) {
  double ea = e_min + i*de;
  double eb = ea + de;

  // Integracion independiente
  double fa = wave_opt(ea,dx,&grid);
  double fb = wave_opt(eb,dx,&grid);

  if (fa*fb < 0.0) {
    double emid = refine(ea,eb,fa);
    update_global(emid); 
  }
}
    \end{lstlisting}
    \end{block}
  \end{column}

  % Columna Derecha: Sección Crítica
  \begin{column}{0.48\textwidth}
    \begin{block}{Ordenamiento (Sección Crítica)}
    \begin{lstlisting}[language=C,style=cppBeamer,basicstyle=\tiny\ttfamily]
#pragma omp critical
{ // Los niveles E_n deben estar ordenados
  if (found < MAX_LEVELS) {
    int pos = found;
    // Insercion por desplazamiento
    while(pos>0 && en[pos-1]>emid) {
      en[pos] = en[pos-1];
      pos--;
    }
    en[pos] = emid;
    found++;
  }
}
    \end{lstlisting}
    \end{block}
  \end{column}

\end{columns}

\vspace{-0.4cm}

% --- Parte Inferior: Comentario de Éxito ---
\begin{successbox}{Estrategia de Carga: \texttt{schedule(dynamic)}}
La bisección de algunas raíces requiere más iteraciones que otras dependiendo de la forma del potencial. El esquema \texttt{dynamic} asegura que los hilos que terminan rápido tomen nuevas tareas del "pool" de intervalos, manteniendo el uso de CPU cercano al $100\%$ y evitando hilos ociosos.
\end{successbox}

\end{frame}

% ─────────────────────────────────────────────────────────────
\begin{frame}[fragile]{OMP: Sinc-DVR — Balanceo de Carga y Vectorización SIMD}

% --- Parte Superior: Dos Columnas de Código ---
\begin{columns}[T]

  % Columna Izquierda: Balanceo Triangular
  \begin{column}{0.48\textwidth}
    \begin{block}{Balanceo en Triángulo Superior}
    \begin{lstlisting}[language=C,style=cppBeamer,basicstyle=\tiny\ttfamily]
// Guided: compensa desbalance j=i+1..n
#pragma omp parallel for schedule(guided)
for (int i=0; i<n; i++) {
  H[i*n+i] = diag_kin + v_i; 

  for (int j=i+1; j<n; j++) {
    const int diff = i-j;
    // Bit-AND: paridad en 1 ciclo
    const double sign = (diff&1) ? -1.0 : 1.0;
    H[i*n+j] = (sign*dx2_inv)/(double)(diff*diff);
  }
}
    \end{lstlisting}
    \end{block}
  \end{column}

  % Columna Derecha: Vectorización
  \begin{column}{0.48\textwidth}
    \begin{block}{Vectorización del Kernel (AVX)}
    \begin{lstlisting}[language=C,style=cppBeamer,basicstyle=\tiny\ttfamily]
// -ffast-math + aligned_alloc(64,...) 
// Genera instrucciones AVX-512
// vmovapd: carga 8 doubles en 1 ciclo
for (int j=i+1; j<n; j++) {
  double inv_d2=1.0/(double)((i-j)*(i-j));
  double sign=((i-j)&1)?-1.0:1.0;
  H[i*n+j]=(sign*dx2_inv)*inv_d2;
}
    \end{lstlisting}
    \end{block}
    \vspace{-0.2cm}
    
  \end{column}

\end{columns}

\vspace{0.3cm}

% --- Parte Inferior: Explicación del Scheduling ---
\begin{conclusionbox}{Justificación de \texttt{schedule(guided)}}
Para $i=0$, el bucle interno realiza $N-1$ iteraciones, mientras que para $i=N-1$ realiza cero. \texttt{Guided} asigna bloques de iteraciones decrecientes: los hilos reciben mucho trabajo al inicio y "pedazos" más pequeños al final para asegurar que todos terminen al mismo tiempo.
\end{conclusionbox}

\end{frame}

% ============================================================
\section{Resultados y Escalabilidad}
% ============================================================

\begin{frame}{Ley de Amdahl y Speedup por Método ($p = 1\ldots 8$ hilos)}
\begin{columns}[T]
\column{0.52\textwidth}
\begin{center}
\begin{tikzpicture}
  \begin{axis}[
    width=6.8cm, height=4.8cm,
    xlabel={Número de hilos $p$},
    ylabel={Speedup $S = T_1/T_p$},
    xmin=1, xmax=8,
    ymin=0.5, ymax=10,
    grid=major, grid style={gray!15},
    tick label style={font=\small},
    label style={font=\small},
    xtick={1,2,4,6,8},
    legend style={font=\tiny, at={(0.03,0.97)}, anchor=north west},
    legend columns=1
  ]
    % Ideal
    \addplot[black, dashed, domain=1:8, samples=8] {x};
    % Shooting (super-linear)
    \addplot[UDsecond, very thick, mark=square*,mark size=3]
      coordinates{(1,1)(2,2.0)(4,4.1)(6,6.2)(8,8.27)};
    % FDM
    \addplot[UDprimary, very thick, mark=triangle*,mark size=3]
      coordinates{(1,1)(2,1.01)(4,1.01)(6,1.02)(8,1.02)};
    % Base SHO
    \addplot[UDsuccess, very thick, mark=*,mark size=3]
      coordinates{(1,1)(2,1.2)(4,1.5)(6,1.7)(8,1.84)};
    % Sinc-DVR
    \addplot[UDaccent!80!black, very thick, mark=diamond*,mark size=3]
      coordinates{(1,1)(2,1.0)(4,0.99)(6,0.99)(8,0.99)};
    \addlegendentry{Ideal lineal}
    \addlegendentry{Shooting $S=8.27$}
    \addlegendentry{FDM $S=1.02$}
    \addlegendentry{Base SHO $S=1.84$}
    \addlegendentry{Sinc-DVR $S=0.99$}
  \end{axis}
\end{tikzpicture}
\end{center}

\column{0.44\textwidth}
\begin{block}{Tabla de Resultados ($p=8$ hilos)}
\centering\small
\begin{tabular}{lcccc}
\toprule
\textbf{Método} & $T_{ser}$ & $T_{opt}$ & $T_{omp8}$ & $S$\\
\midrule
Shooting & 0.232 & 0.031 & 0.028 & \textbf{8.27}\\
FDM      & 7.191 & 0.038 & 7.018 & 1.02\\
Base SHO & 0.179 & 0.106 & 0.097 & 1.84\\
Sinc-DVR & 0.110 & 0.110 & 0.110 & 0.99\\
\bottomrule
\end{tabular}
\end{block}
\vspace{0.135cm}
\begin{alertbox}{Speedup Superlineal en Shooting}
$E_p > 1.0$ porque al dividir el dominio de energía entre 8 hilos, cada hilo integra un sub-intervalo cuya malla de potencial \textbf{cabe íntegramente en la caché L1} del núcleo, eliminando fallos de caché.
\end{alertbox}
\end{columns}
\end{frame}

% ─────────────────────────────────────────────────────────────
\begin{frame}{Eficiencia Paralela y Ley de Amdahl: Discusión}
\begin{columns}[T]
\column{0.50\textwidth}
\begin{center}
\begin{tikzpicture}
  \begin{axis}[
    width=6.4cm, height=4.4cm,
    title={\small Eficiencia Paralela $E_p = S/p$},
    xlabel={Número de hilos $p$},
    ylabel={Eficiencia $E_p$},
    xmin=1, xmax=8,
    ymin=0, ymax=1.2,
    grid=major, grid style={gray!15},
    tick label style={font=\small},
    label style={font=\small},
    legend style={font=\tiny, at={(0.97,0.97)}, anchor=north east},
    xtick={1,2,4,6,8},
    yticklabel style={/pgf/number format/fixed,/pgf/number format/precision=2}
  ]
    \addplot[black,dashed,domain=1:8]{1.0};
    \addplot[UDsecond,very thick,mark=square*,mark size=2.5]
      coordinates{(1,1.0)(2,1.0)(4,1.03)(6,1.03)(8,1.03)};
    \addplot[UDsuccess,very thick,mark=*,mark size=2.5]
      coordinates{(1,1.0)(2,0.60)(4,0.38)(6,0.28)(8,0.23)};
    \addplot[UDprimary,very thick,mark=triangle*,mark size=2.5]
      coordinates{(1,1.0)(2,0.51)(4,0.26)(6,0.17)(8,0.13)};
    \addlegendentry{Ideal}
    \addlegendentry{Shooting}
    \addlegendentry{Base SHO}
    \addlegendentry{FDM}
  \end{axis}
\end{tikzpicture}
\end{center}

\column{0.46\textwidth}
\begin{block}{4 Fenómenos Clave de Rendimiento}
\begin{enumerate}\setlength\itemsep{3pt}
  \item \textbf{Speedup superlineal (Shooting):} División del dominio $\rightarrow$ datos en L1/L2 por núcleo $\rightarrow$ cero fallos de caché $\rightarrow$ $E_p > 1$

  \item \textbf{Ley de Amdahl en FDM:} La diagonalización densa (fracción serial $f \approx 0.99$) fija $S_{max} \approx 1.02$, sin importar cuántos hilos se añadan

  \item \textbf{Saturación en Sinc-DVR:} El costo de comunicación y sincronización entre hilos para $N=500$ supera las ganancias aritméticas

  \item \textbf{Eficiencia decreciente en SHO:} Para $N=400$, el overhead de gestión de hilos comienza a competir con el trabajo útil
\end{enumerate}
\end{block}
\end{columns}
\end{frame}

% ─────────────────────────────────────────────────────────────
\begin{frame}{Escalabilidad en Escala Log-Log: Comportamiento Asintótico}
\begin{columns}[T]
\column{0.50\textwidth}
\begin{center}
\begin{tikzpicture}
  \begin{axis}[
    width=5.6cm, height=3.6cm,
    title={\small Shooting: Escalabilidad Log-Log},
    xlabel={$N_{steps}$}, ylabel={Tiempo (s)},
    xmode=log, ymode=log,
    xmin=100, xmax=3000,
    ymin=1e-4, ymax=1,
    grid=major,grid style={gray!15},
    tick label style={font=\tiny},
    label style={font=\small},
    legend style={font=\tiny,at={(0.03,0.97)},anchor=north west}
  ]
    \addplot[gray!60,very thick] coordinates {(100,2e-1) (3000,2e-1)};
    \addplot[UDsecond,thick,mark=*,mark size=1.5]
      coordinates{(100,8e-2)(500,0.12)(1000,0.22)(2000,0.23)};
    \addplot[UDsecond!70,thick,mark=square*,mark size=1.5]
      coordinates{(100,4e-2)(500,6e-2)(1000,1.1e-1)(2000,1.15e-1)};
    \addplot[UDsecond!40,thick,mark=triangle*,mark size=1.5]
      coordinates{(100,1.5e-2)(500,2.5e-2)(1000,4e-2)(2000,5e-2)};
    \addplot[UDsecond!20,thick,mark=diamond*,mark size=1.5]
      coordinates{(100,1e-2)(500,1.5e-2)(1000,2.5e-2)(2000,2.8e-2)};
    \addlegendentry{Serial}
    \addlegendentry{2 hilos}
    \addlegendentry{4 hilos}
    \addlegendentry{8 hilos}
  \end{axis}
\end{tikzpicture}
\end{center}
\begin{successbox}{Pendientes Paralelas = Speedup Lineal}
Líneas equidistantes en log-log $\Rightarrow$ el factor de escala (speedup) es constante con $N$: paralelismo verdaderamente escalable.
\end{successbox}

\column{0.46\textwidth}
\begin{center}
\begin{tikzpicture}
  \begin{axis}[
    width=5.6cm, height=3.6cm,
    title={\small FDM: Techo de Amdahl},
    xlabel={$N$}, ylabel={Tiempo (s)},
    xmode=log, ymode=log,
    xmin=100, xmax=2000,
    ymin=1e-3, ymax=20,
    grid=major,grid style={gray!15},
    tick label style={font=\tiny},
    label style={font=\small},
    legend style={font=\tiny,at={(0.03,0.97)},anchor=north west}
  ]
    \addplot[UDprimary,thick,mark=*,mark size=1.5]
      coordinates{(100,0.05)(200,0.4)(500,3.5)(1000,7.2)(2000,15)};
    \addplot[UDprimary!50,thick,mark=square*,mark size=1.5]
      coordinates{(100,0.045)(200,0.37)(500,3.4)(1000,7.1)(2000,14.8)};
    \addplot[UDprimary!20,thick,mark=triangle*,mark size=1.5]
      coordinates{(100,0.04)(200,0.36)(500,3.35)(1000,7.0)(2000,14.7)};
    \addlegendentry{Serial}
    \addlegendentry{4 hilos}
    \addlegendentry{8 hilos}
  \end{axis}
\end{tikzpicture}
\end{center}
\begin{alertbox}{Curvas Convergentes = Ley de Amdahl}
Para $N$ grande, las curvas convergen: la diagonalización $O(N^3)$ de LAPACK domina y no se paraleliza.
\end{alertbox}
\end{columns}
\end{frame}

% ============================================================
\section{Comparación Final}
% ============================================================

% --- DIAPOSITIVA 1: Tabla de Métricas Comparativas ---
\begin{frame}{Comparación Final Integral: Métricas de Rendimiento}

\begin{center}
\small % Ajuste para que la tabla sea la protagonista
\begin{tabular}{lp{1.8cm}p{1.8cm}ccp{2.2cm}}
\toprule
\textbf{Método} & \textbf{Error $E_{0..3}$} & \textbf{Convergencia} & \textbf{$S_{serial}$} & \textbf{$S_{OMP8}$} & \textbf{Escalabilidad}\\
\midrule
\textbf{Shooting} & $10^{-5}$ & Lineal $O(N)$ & $3.5\times$ & $\mathbf{8.27\times}$ & Excelente \\
\midrule
\textbf{FDM} & $10^{-4}$ & Algebraica $O(h^2)$ & $>100\times$ & $1.02\times$ & Nula (LAPACK) \\
\midrule
\textbf{Base SHO} & $10^{-6}$ & Espectral & $2.5\times$ & $1.84\times$ & Moderada \\
\midrule
\textbf{Sinc-DVR} & $10^{-7}..10^{-8}$ & Exponencial & $2.2\times$ & $0.99\times$ & Nula (Memoria) \\
\bottomrule
\end{tabular}
\end{center}

\vspace{0.6cm}
\begin{block}{Notas sobre los Speedups ($S$)}
\footnotesize
\begin{itemize}
    \item $S_{serial}$: Mejora respecto al código base sin optimizar.
    \item $S_{OMP8}$: Escalabilidad en 8 hilos (Shooting presenta super-linearidad por caché).
\end{itemize}
\end{block}

\end{frame}


% --- DIAPOSITIVA 2: Recomendaciones y Casos de Uso ---
\begin{frame}{Conclusiones: ¿Qué método utilizar?}

\vspace{0.5cm}
\begin{columns}[T]

% Columna 1: Escalabilidad
\column{0.32\textwidth}
\begin{successbox}{\faAward\ Escalabilidad}
\textbf{Shooting Method}\\
\vspace{0.2cm}
\small
Ideal para escenarios HPC masivos. Su paralelismo de tareas permite una eficiencia casi perfecta.
\end{successbox}

% Columna 2: Precisión
\column{0.32\textwidth}
\begin{physbox}{\faBullseye\ Precisión}
\textbf{Sinc-DVR}\\
\vspace{0.2cm}
\small
Recomendado para cálculos de referencia donde el error debe ser $< 10^{-7}$ con un $N$ mínimo.
\end{physbox}

% Columna 3: Balance
\column{0.32\textwidth}
\begin{conclusionbox}{\faBalanceScale\ Balance}
\textbf{Base SHO}\\
\vspace{0.2cm}
\small
El "punto dulce": buena precisión ($10^{-6}$) con una escalabilidad aceptable en OpenMP.
\end{conclusionbox}

\end{columns}

\vspace{0.8cm}
\begin{alertbox}{Advertencia sobre FDM}
\footnotesize
Aunque el speedup serial es masivo ($>100\times$), su dependencia de la diagonalización de matrices grandes lo hace el menos apto para paralelismo de memoria compartida tradicional.
\end{alertbox}

\end{frame}

% ─────────────────────────────────────────────────────────────
\begin{frame}{Precisión de los Primeros 4 Niveles: Error vs. Analítico}
\begin{center}
\begin{tikzpicture}
  \begin{axis}[
    ybar, bar width=9pt,
    width=11.2cm, height=4.8cm,
    symbolic x coords={$E_0$,$E_1$,$E_2$,$E_3$},
    xtick=data,
    ylabel={Error relativo $|\Delta E/E_{analítico}|$},
    ymode=log,
    ymin=1e-9, ymax=1e-2,
    axis lines=left,
    grid=major, grid style={gray!15},
    tick label style={font=\small},
    label style={font=\small},
    legend style={font=\small, at={(0.98,0.98)}, anchor=north east, fill=white},
    legend columns=2
  ]
    \addplot[fill=UDsecond!70, draw=UDsecond]
      coordinates{($E_0$,3e-5)($E_1$,4e-5)($E_2$,5e-5)($E_3$,7e-5)};
    \addplot[fill=UDprimary!70, draw=UDprimary]
      coordinates{($E_0$,5e-4)($E_1$,6e-4)($E_2$,8e-4)($E_3$,1.2e-3)};
    \addplot[fill=UDsuccess!70, draw=UDsuccess]
      coordinates{($E_0$,3e-6)($E_1$,4e-6)($E_2$,5e-6)($E_3$,8e-6)};
    \addplot[fill=UDaccent!80, draw=UDaccent!80!black]
      coordinates{($E_0$,5e-8)($E_1$,6e-8)($E_2$,8e-8)($E_3$,1.2e-7)};
    \addlegendentry{Shooting (N=1000)}
    \addlegendentry{FDM (N=1000)}
    \addlegendentry{Base SHO (N=400)}
    \addlegendentry{Sinc-DVR (N=200)}
  \end{axis}
\end{tikzpicture}
\end{center}
\begin{center}
\small
\textbf{Anarmonicidad observada}: el error crece con $n$ en todos los métodos, reflejando que los estados excitados exploran regiones del potencial más alejadas del equilibrio.
\end{center}
\end{frame}

% ─────────────────────────────────────────────────────────────
\begin{frame}{Radar de Rendimiento: Comparación Multidimensional}
\begin{center}
\begin{tikzpicture}[scale=0.88]
  % Simple radar chart as polygon overlay on axes
  \def\R{2.8}
  % Draw axes
  \foreach \angle/\label in {90/Precisión, 162/Escalabilidad, 234/Opt.Serial, 306/Costo Mem., 18/Vel. Constr.} {
    \draw[-,gray!50] (0,0) -- (\angle:\R);
    \node[font=\tiny\bfseries] at (\angle:\R+0.35) {\label};
  }
  % Concentric circles
  \foreach \r in {0.7,1.4,2.1,2.8} {
    \draw[gray!30] (0,0) circle(\r);
  }
  % Shooting: Prec=0.6, Escal=1.0, Opt=0.7, Mem=1.0, Vel=0.9
  \draw[UDsecond,very thick,fill=UDsecond!20,opacity=0.7]
    (90:0.6*\R) -- (162:1.0*\R) -- (234:0.7*\R) -- (306:1.0*\R) -- (18:0.9*\R) -- cycle;
  % FDM: Prec=0.5, Escal=0.1, Opt=1.0, Mem=0.6, Vel=1.0
  \draw[UDprimary,very thick,fill=UDprimary!20,opacity=0.7]
    (90:0.5*\R) -- (162:0.1*\R) -- (234:1.0*\R) -- (306:0.6*\R) -- (18:1.0*\R) -- cycle;
  % Base SHO: Prec=0.8, Escal=0.5, Opt=0.6, Mem=0.7, Vel=0.7
  \draw[UDsuccess,very thick,fill=UDsuccess!20,opacity=0.7]
    (90:0.8*\R) -- (162:0.5*\R) -- (234:0.6*\R) -- (306:0.7*\R) -- (18:0.7*\R) -- cycle;
  % Sinc-DVR: Prec=1.0, Escal=0.1, Opt=0.5, Mem=0.3, Vel=0.6
  \draw[UDaccent!80!black,very thick,fill=UDaccent!20,opacity=0.7]
    (90:1.0*\R) -- (162:0.1*\R) -- (234:0.5*\R) -- (306:0.3*\R) -- (18:0.6*\R) -- cycle;
  % Legend
  \node[draw,fill=white,anchor=north west,font=\tiny] at (2.5,1.8) {
    \begin{tabular}{ll}
      \textcolor{UDsecond}{$\blacksquare$} & Shooting\\
      \textcolor{UDprimary}{$\blacksquare$} & FDM\\
      \textcolor{UDsuccess}{$\blacksquare$} & Base SHO\\
      \textcolor{UDaccent!80!black}{$\blacksquare$} & Sinc-DVR
    \end{tabular}
  };
\end{tikzpicture}
\end{center}
\end{frame}

% ============================================================
\section{Conclusiones}
% ============================================================

\begin{frame}{Conclusiones}
\begin{columns}[T]
\column{0.50\textwidth}
\begin{successbox}{\faCheck\ 1. Primacía del Algoritmo sobre el Hardware}
La mejora de $\mathbf{>100\times}$ en FDM demuestra que la eficiencia es ante todo una cuestión de \textbf{diseño estructural}: pre-cómputo, localidad de datos y explotación de la esparcidad superan cualquier ganancia de núcleos adicionales.
\end{successbox}
\vspace{0.27cm}
\begin{successbox}{\faCheck\ 2. Speedup Superlineal y Jerarquía de Memoria}
El fenómeno $E_p > 1.0$ en \textbf{Shooting} no contradice la teoría: al dividir el dominio de energías entre núcleos, los datos residen en \textbf{caché L1/L2 por núcleo}, eliminando latencia de RAM.
\end{successbox}
\vspace{0.27cm}
\begin{alertbox}{\faExclamationTriangle\ 3. Ley de Amdahl — Cuellos Secuenciales}
FDM y Sinc-DVR exhiben $S \approx 1$: la diagonalización densa de LAPACK ($O(N^3)$, serial) es la fracción no paralelizable que limita el speedup máximo, sin importar cuántos hilos se añadan.
\end{alertbox}

\column{0.46\textwidth}
\begin{physbox}{\faBook\ 4. Selección de Método por Escenario}
\begin{enumerate}\setlength\itemsep{4pt}
  \item \textbf{Escalabilidad masiva} (HPC multihilo):\\
  \hspace{1em}$\Rightarrow$ \textbf{Shooting} — paralelismo de tareas, $E_p>1$
  \item \textbf{Precisión extrema} ($N$ reducido):\\
  \hspace{1em}$\Rightarrow$ \textbf{Sinc-DVR} — convergencia exponencial, $N<300$
  \item \textbf{Balance precisión/costo}:\\
  \hspace{1em}$\Rightarrow$ \textbf{Base SHO} — espectral moderado + OpenMP razonable
  \item \textbf{Mallas densas con optimización serial}:\\
  \hspace{1em}$\Rightarrow$ \textbf{FDM} — estructura esparsa explotada
\end{enumerate}
\end{physbox}
\vspace{0.225cm}
\begin{conclusionbox}{\faStar\ Síntesis Final}
El éxito en computación cuántica numérica depende de \textbf{alinear el modelo físico con la arquitectura del procesador}: optimizar no solo el conteo de operaciones, sino el \textbf{flujo de datos a través de la jerarquía de memoria L1/L2/L3/RAM}.
\end{conclusionbox}
\end{columns}
\end{frame}

% ─── DIAPOSITIVA FINAL ───────────────────────────────────────
{
\setbeamertemplate{headline}{}
\setbeamertemplate{footline}{}
\begin{frame}
\begin{tikzpicture}[remember picture, overlay]
  \fill[UDdark] (current page.south west) rectangle (current page.north east);
  \fill[UDprimary!70,opacity=0.7]
    ([yshift=-0.35\paperheight]current page.north west) rectangle
    ([yshift=0.35\paperheight]current page.south east);
  \fill[UDsecond]
    ([xshift=0pt]current page.south west) rectangle
    ([yshift=0.018\paperheight,xshift=\paperwidth]current page.south west);
\end{tikzpicture}
\begin{center}
\vspace{1.35cm}
{\LARGE\bfseries\color{white} ¡Gracias!}\\[0.6cm]
{\color{UDlight}\large Optimización y Paralelización del Potencial de Morse}\\[0.3cm]
{\color{UDlight}\small Primer Parcial — Sistemas Distribuidos 2026-1}\\[0.6cm]
\begin{tcolorbox}[colback=UDprimary!20,colframe=UDaccent,width=9cm,
  fonttitle=\bfseries\small\color{white},title={Autores}]
\centering\color{white}
\textbf{Camilo Huertas} \quad | \quad \textbf{Sebastián Rodriguez}\\[2pt]
{\small Programa Académico de Física\\
Universidad Distrital Francisco José de Caldas}\\[2pt]
{\small Docente: Carlos Andrés Gómez Vasco}
\end{tcolorbox}
\vspace{0.36cm}
{\color{UDaccent}\small 27 de marzo de 2026}
\end{center}
\end{frame}
}

\end{document}